\subsection{COMPUTATIONAL COST}
\label{app:computational_cost}

Adaptive inference trades the cost of constructing a sampling density for
a smaller operator evaluation. Table~\ref{tab:inference_time_sampled}
quantifies this trade-off across all eight benchmarks at the three
evaluation budgets. All methods use the same frozen backbone. Uniform
provides a reference at matched point counts, while full mesh measures the
cost of evaluating all reference points. The reported latency includes
online refinement and prediction at the selected points, with reusable
preprocessing performed in advance. Reconstruction onto the full mesh is
excluded.

\begin{table}[!ht]
\centering
\begingroup
\footnotesize
\setlength{\tabcolsep}{3.0pt}
\renewcommand{\arraystretch}{1.08}
\begin{tabular}{llrrrrr}
\toprule
Task & Budget & Full mesh & Uniform & PI & CPG & CA \\
\midrule
Airfoil2D & 37.5\% & $5.280\!\pm\!0.009$ & $2.438\!\pm\!0.018$ & $2.539\!\pm\!0.016$ & $4.435\!\pm\!0.039$ & $3.110\!\pm\!0.031$ \\
 & 50\% &  & $3.037\!\pm\!0.005$ & $3.176\!\pm\!0.006$ & $5.050\!\pm\!0.012$ & $3.775\!\pm\!0.012$ \\
 & 75\% &  & $4.049\!\pm\!0.009$ & $4.218\!\pm\!0.013$ & $6.104\!\pm\!0.018$ & $4.847\!\pm\!0.016$ \\
\addlinespace[2pt]
Pipe & 37.5\% & $7.407\!\pm\!0.004$ & $3.238\!\pm\!0.002$ & $3.382\!\pm\!0.004$ & $5.067\!\pm\!0.005$ & $4.095\!\pm\!0.005$ \\
 & 50\% &  & $4.038\!\pm\!0.005$ & $4.199\!\pm\!0.005$ & $5.891\!\pm\!0.007$ & $4.923\!\pm\!0.007$ \\
 & 75\% &  & $5.780\!\pm\!0.021$ & $6.025\!\pm\!0.016$ & $7.721\!\pm\!0.020$ & $6.774\!\pm\!0.028$ \\
\addlinespace[2pt]
Darcy & 37.5\% & $3.525\!\pm\!0.001$ & $1.929\!\pm\!0.002$ & $1.994\!\pm\!0.003$ & $3.469\!\pm\!0.007$ & $2.675\!\pm\!0.004$ \\
 & 50\% &  & $2.212\!\pm\!0.001$ & $2.301\!\pm\!0.003$ & $3.785\!\pm\!0.017$ & $3.002\!\pm\!0.015$ \\
 & 75\% &  & $2.907\!\pm\!0.005$ & $3.030\!\pm\!0.004$ & $4.536\!\pm\!0.014$ & $3.738\!\pm\!0.006$ \\
\addlinespace[2pt]
Elasticity$\dagger$ & 37.5\% & $1.264\!\pm\!0.013$ & $0.987\!\pm\!0.007$ & $1.009\!\pm\!0.006$ & $2.086\!\pm\!0.013$ & $1.524\!\pm\!0.024$ \\
 & 50\% &  & $1.027\!\pm\!0.004$ & $1.047\!\pm\!0.003$ & $2.209\!\pm\!0.003$ & $1.598\!\pm\!0.017$ \\
 & 75\% &  & $1.166\!\pm\!0.003$ & $1.184\!\pm\!0.004$ & $2.401\!\pm\!0.038$ & $1.756\!\pm\!0.012$ \\
\addlinespace[2pt]
Plasticity & 37.5\% & $2.049\!\pm\!0.006$ & $1.259\!\pm\!0.002$ & $1.289\!\pm\!0.003$ & $2.516\!\pm\!0.005$ & $1.657\!\pm\!0.003$ \\
 & 50\% &  & $1.430\!\pm\!0.002$ & $1.470\!\pm\!0.002$ & $2.716\!\pm\!0.007$ & $1.833\!\pm\!0.004$ \\
 & 75\% &  & $1.708\!\pm\!0.062$ & $1.772\!\pm\!0.071$ & $3.049\!\pm\!0.065$ & $2.152\!\pm\!0.078$ \\
\addlinespace[2pt]
ShapeNet-Car & 37.5\% & $16.359\!\pm\!0.011$ & $6.856\!\pm\!0.004$ & $6.976\!\pm\!0.005$ & $10.609\!\pm\!0.026$ & $8.401\!\pm\!0.027$ \\
 & 50\% &  & $8.644\!\pm\!0.008$ & $8.773\!\pm\!0.006$ & $12.404\!\pm\!0.025$ & $10.259\!\pm\!0.027$ \\
 & 75\% &  & $12.388\!\pm\!0.004$ & $12.618\!\pm\!0.006$ & $16.258\!\pm\!0.029$ & $14.148\!\pm\!0.029$ \\
\addlinespace[2pt]
DrivAerML & 37.5\% & $4601.5\!\pm\!140.1$ & $1726.7\!\pm\!46.4$ & $1758.6\!\pm\!47.6$ & $2228.2\!\pm\!61.2$ & $1765.0\!\pm\!46.6$ \\
 & 50\% &  & $2306.9\!\pm\!62.1$ & $2347.1\!\pm\!62.2$ & $2825.2\!\pm\!74.3$ & $2345.7\!\pm\!59.6$ \\
 & 75\% &  & $3460.8\!\pm\!97.0$ & $3520.6\!\pm\!97.2$ & $4005.8\!\pm\!109.6$ & $3505.5\!\pm\!94.4$ \\
\addlinespace[2pt]
SuperWing & 37.5\% & $13.602\!\pm\!0.010$ & $5.671\!\pm\!0.002$ & $5.909\!\pm\!0.003$ & $8.610\!\pm\!0.006$ & $6.659\!\pm\!0.007$ \\
 & 50\% &  & $7.401\!\pm\!0.005$ & $7.744\!\pm\!0.005$ & $10.447\!\pm\!0.007$ & $8.487\!\pm\!0.006$ \\
 & 75\% &  & $10.483\!\pm\!0.010$ & $10.987\!\pm\!0.008$ & $13.690\!\pm\!0.010$ & $11.763\!\pm\!0.010$ \\
\bottomrule
\end{tabular}
\endgroup
\caption{\textbf{Prepared sampled-field request latency (ms).}
Values are means $\pm$ sample standard deviations across three cases and three sampling seeds.
The full-mesh column uses all reference points and three cases, and is shown once per task.
$\dagger$ The Elasticity timing configuration uses a $25\%$ coarse set and a two-block head.
Complete-field reconstruction is measured separately in Table~\ref{tab:inference_time_reconstruction}.}
\label{tab:inference_time_sampled}
\end{table}


\paragraph{Refinement overhead.}
PI remains close to Uniform, adding $1.5$--$4.8\%$ latency across the
24 task--budget pairs. Its geometry-based density can be reused across
predictions. CPG incurs a larger overhead because it evaluates a complete
coarse prediction before selecting additional points. CA obtains its
refinement signal from early features and reuses cached pointwise
quantities during the final evaluation. It reduces latency by
$12.3$--$34.1\%$ relative to CPG, with a reduction in every measured
setting. Thus, the learned refinement signal lowers the cost of
solution-aware sampling while retaining its ability to adapt to the input.

\paragraph{Realized speedup.}
At a $37.5\%$ point budget, CA is faster than full-mesh inference on seven
of the eight benchmarks. The speedups reach $1.95\times$ on ShapeNet-Car,
$2.61\times$ on DrivAerML, and $2.04\times$ on SuperWing. These gains are
more pronounced when evaluating the spatial points accounts for a large
share of the computation. On Elasticity, CA takes $1.52$~ms compared with
$1.26$~ms for full mesh: refinement overhead outweighs the savings from
evaluating fewer points. Point reduction therefore does not imply a
proportional reduction in latency. The realized gain depends on how much
operator computation is removed relative to the cost of acquiring the
refinement signal.
